% Imporditud: tools/import_eksperimendid.py
% Allikas: Eesti füüsikaolümpiaadi eksperimendi ülesannete
%   kogu (Rael Kalda, 2023), ülesanne Ü16, lahendus L16.
\setAuthor{EFO žürii}
\setRound{piirkonnavoor}
\setYear{2008}
\setNumber{P E1}
\setDifficulty{1}
\setTopic{vedelike mehaanika}
\setVahendid{mutter, polt, mõõtejoonlaud, niit, anum veega}
\setPunktid{10}
\setAllikas{Eksperimendi kogu Ü16, lahendus lk 39}
\setLahendusseis{toores}

\prob{Mutter}
Mutter ja polt on valmistatud samast ainest. Kumma mass on suurem ja mitu korda?

\hint

\solu
Laseme mutri anumasse, mõõdame veetaseme tõusu. Sama teeme poldiga. Mõõdame pudeli diameetri. Arvutame veetaseme kõrguste muutuste suhte, mis võrdub masside suhtega. m = $\rho$V ja V = Sp$\Delta$h, seega:

mp = Vp = $\Delta$hp . mm Vm $\Delta$hm

Hindamine: Teoreetiline mudel (5 p). Mõõtmiste teostamine ja suhte arvutamine (3 p); tulemus on tõepärane (1 p). Mõõtmisi on korratud mitu korda (1 p).
\probend
